\documentclass{beamer}
\usepackage[utf8]{inputenc}
\usepackage[spanish]{babel}
\usetheme{Madrid}
\title{Estadística Descriptiva}
\subtitle{Tema 1: medidas de tendencia central y dispersión}
\author{Ángel Berihuete}
\date{Curso 2026--2027}

\begin{document}

\begin{frame}
  \titlepage
\end{frame}

\begin{frame}{Media y varianza}
  La media de una muestra $x_1,\dots,x_n$ es
  \[
    \bar{x} = \frac{1}{n}\sum_{i=1}^{n} x_i .
  \]
  La varianza muestral es
  \[
    s^2 = \frac{1}{n-1}\sum_{i=1}^{n}(x_i-\bar{x})^2 .
  \]
\end{frame}

\begin{frame}{Histograma}
  El histograma de los tiempos de espera del géiser \texttt{faithful}:
  % Aquí iría una figura generada con R
\end{frame}

\begin{frame}{Medidas de posición}
  \begin{itemize}
    \item Media
    \item Mediana
    \item Cuartiles
  \end{itemize}
\end{frame}

\end{document}
